\section{The Present Situation, Stated as the Acts State It}

What follows describes the legal position created by the acts set out
above, as far as possible in their own terms, as of 25 July 2026. It is
not legal advice, not a canonical opinion, and not a judgment about any
person. Where the acts leave a question open, this section says so and
stops.

\subsection{Canonical status}

The Priestly Fraternity of Saint Peter is a clerical society of apostolic
life of pontifical right. It was erected as such by decree of the
Pontifical Commission \work{Ecclesia Dei} on 18 October 1988,
Prot.\ N.~234/88, in exercise of faculty 3~a) of the papal rescript
\latin{Quia peculiare munus} of the same day, and its constitutions were
definitively approved on 29 June 2003, Prot.\ N.~108/2003. No act located
for this account has suppressed, amalgamated, reconstituted or amended
that erection.

Its members take no religious vows. Its clerics are incardinated in the
society itself under canon 736 §\,1, save as its own law provides
otherwise. In matters of internal governance and discipline it is subject
immediately and exclusively to the Apostolic See under canon 593 as
applied by canon 732. In public worship, the care of souls and other works
of the apostolate its members are subject to the diocesan bishop under
canon 738 §\,2, with attention to canons 679--683, which the erection
decree names expressly. Its houses require the previous written consent of
the diocesan bishop under canon 733 §\,1 and carry, once erected, the
rights of canon 611 which the erection decree invokes.

\subsection{Supervising authority}

The dicastery competent for the Fraternity is the Dicastery for
Institutes of Consecrated Life and Societies of Apostolic Life. The chain
is: faculty 6 of the 1988 rescript, granting the Pontifical Commission
the exercise of the Holy See's authority \latin{donec aliter provideatur};
the motu proprio of 17 January 2019 decommissioning the Commission and
assigning its tasks to the Congregation for the Doctrine of the Faith,
expressly citing that faculty; and \work{Traditionis custodes} art.~6 of
16 July 2021, placing institutes and societies erected by the Commission
under the dicastery for consecrated life. The Fraternity dates its own
change of supervisor to 2021.

\subsection{Liturgical position}

\begin{itemize}
\item By the decree of Francis of 11 February 2022, each and every member
holds the faculty to celebrate Mass, the sacraments and other sacred
rites and to fulfil the Divine Office according to the typical editions
of the Missal, Ritual, Pontifical and Roman Breviary in force in 1962.
\item That faculty operates by itself in the Fraternity's own churches
and oratories; elsewhere it may be used only with the consent of the
local Ordinary, private Masses excepted.
\item The decree adds, the foregoing being duly observed, that the Holy
Father ``also advises that careful thought be given, as far as can be
done,'' to the provisions of \work{Traditionis custodes}. That clause is
hortatory in its verb and is not stated as a condition of the grant.
\item The general law remains \work{Traditionis custodes}: the reformed
books are the unique expression of the \latin{lex orandi} of the Roman
Rite (art.~1); authorizing the 1962 Missal in a diocese is the exclusive
competence of its bishop according to the guidelines of the Apostolic See
(art.~2); the use of a parish church, the erection of a personal parish,
and permission for priests ordained after 16 July 2021 are dispensations
reserved in a special way to the Apostolic See (rescript of 20 February
2023); and a diocesan bishop may permit the 1952 Ritual only to
canonically erected personal parishes and may not permit the pre-conciliar
Pontifical (\work{Responsa} of 4 December 2021).
\end{itemize}

\subsection{Government and scale, as the institute reports them}

The superior general is Father John Berg, elected on 9 July 2024 by the
seventh ordinary general chapter for six years, in his third
non-consecutive term. At 1 November 2025 the Fraternity reported 579
members, of whom 394 incardinated: 387 priests (365 incardinated, 16
incorporated \latin{ad annum}, 3 postulants, 3 associated), 30 deacons and
162 non-deacon seminarians, with an average age of 39. Its published
organization chart at the cutoff shows a North American Province, a French
Province, German-speaking and Oceania districts, and two seminaries, at
Wigratzbad in the diocese of Augsburg and at Denton in the diocese of
Lincoln. All of these are self-reported and are recorded here as such.

\subsection{What the acts do not say}

\begin{itemize}
\item They do not say that \work{Traditionis custodes} does not bind the
Fraternity. The proposition is reported by the Fraternity as something
said in a private audience on 4 February 2022; the decree of 11 February
2022 does not state it, and no located act states it.
\item They do not say whether art.~8 of \work{Traditionis custodes}
abrogated the liturgical concessions of 10 September and 18 October 1988.
The 2022 decree re-grants the substance of the first without deciding the
question, and says nothing about the second's extension to guests and to
non-member priests ministering in the Fraternity's churches.
\item They do not resolve the relation between the 2022 faculty and
art.~4 of \work{Traditionis custodes} on priests ordained after 16 July
2021.
\item They do not state a term for the 2022 faculty, provide for its
review, or promulgate it in \work{Acta Apostolicae Sedis}. Absence of a
term is not an assurance of perpetuity.
\item They do not say what the apostolic visitation opened in September
2024 has found. No report, decree or communiqué concluding it had been
published by the Holy See or by the Fraternity at the cutoff.
\item They do not settle whether the exclusive use of the 1962 books may
be required of the Fraternity's own members. The letter of 29 June 2000
says it may not be made the general rule of an institute; the decree of 11
February 2022 grants a faculty and imposes nothing; the published
constitutions require ``faithful observance'' and do not use the word
``exclusive.''
\end{itemize}

\actrecord{The acts in force at 25 July 2026: the erection decree of 18
October 1988 with the rescript \latin{Quia peculiare munus}; the
constitutions definitively approved 29 June 2003; \work{Traditionis
custodes} of 16 July 2021 with the \work{Responsa} of 4 December 2021 and
the rescript of 20 February 2023; and the papal decree of 11 February
2022.}{That the Fraternity is a clerical society of apostolic life of
pontifical right under the dicastery for consecrated life, whose members
hold by papal faculty the use of the four 1962 books in their own
churches and oratories and elsewhere with the local Ordinary's consent,
within a general liturgical law that reserves several dispensations to the
Apostolic See.}{An apostolic visitation is open with no published
outcome; the principal liturgical title has no official publication; and
the relation between that title and several articles of the governing
motu proprio is undetermined on the public record.}
